\section{Discussion}
\label{sec:discussion}

\paragraph{The certificate is a conservative lower bound.}
The empirical coverage under the budget ($0.988$) sits well above the certified floor
($0.75$), and the construction explains why. Theorem~\ref{thm:composed} composes the two
miscoverage events through a union bound, which is tight only when the two failures are
disjoint. In practice the soft-sensor and prognostic failures are \emph{positively}
correlated in the regimes that matter: aggressive operation simultaneously pushes the
soft sensor out of its calibration regime and loads the pump, so the two modules tend to
be stressed together. Positive failure correlation makes the true joint coverage exceed
the Bonferroni bound, and the deliberately conservative Lipschitz constants
($L_1=L_2=0.15$ against a swept $L_1\approx0.024$) widen the gap further. The certificate
should therefore be read as a guarantee --- a floor that is provably met --- not as a
prediction of the achieved coverage. Tightening it is a concrete avenue: a joint
failure-correlation model, via a copula on the two nonconformity scores, multivariate
conformal prediction, or conformal risk control under dependence, would replace the union
bound with a dependence-aware composition and recover much of the margin.

\paragraph{From per-cycle to the operating horizon.}
Theorem~\ref{thm:composed} is a per-cycle, marginal guarantee. A campaign runs over many
cycles, and a naive union of the per-cycle bound over a horizon of $K$ cycles degrades
linearly in $K$. The route to a genuine horizon guarantee is adaptive conformal
inference~\citep{gibbs2021}, which the soft-sensor module already
implements~\citep{busico_m1}: by adjusting the conformal quantile online it secures
long-run coverage under arbitrary, unspecified dependence --- exactly the dependence that
feedback induces across cycles. Lifting the per-cycle certificate to a horizon statement
under this mechanism, in a dynamic rather than quasi-static setting, is the principal
direction of ongoing work and is what connects the present result to a physically
deployed implementation.

\paragraph{The selection term, and when it is not zero.}
The selection penalty $\Delta_{\mathrm{sel},k}$ vanishes here as a property of the loop
(Assumption~\ref{ass:causal}), not as an idealization: the synchronous controller forms
each decision from the previous cycle's information and applies it before the current
measurement is realized. The value of carrying the term explicitly is that it
delineates precisely when conformal validity \emph{would} be lost --- under a non-causal
co-selection in which the optimizer chooses the decision using the same residual that the
coverage check then evaluates. Most treatments of conformal prediction in feedback either
assume the difficulty away or abandon the guarantee; naming and bounding the selection
effect, and showing it is null under realizable timing, is what makes the closed-loop
certificate honest.

\paragraph{Scope and generality.}
The certificate is instantiated for binary-key distillation through a similitude
coordinate built from the relative volatility, the Gilliland reflux coordinate, and the
stripping factor. That coordinate is the unit-specific ingredient; the composition logic
--- per-module conditional coverage, the causal-timing argument, the Bonferroni
composition, and the convex $\psi$-budget --- is unit-agnostic. Extending the framework
to multicomponent separations or to other unit operations requires identifying the
appropriate dimensionless groups for those units, after which the same machinery applies.
A second refinement concerns the Lipschitz constants, here estimated a~priori from a
regime sweep; an online estimator that re-identifies them as the plant ages would let the
budget adapt, and pairs naturally with the adaptive-conformal horizon extension above.
The budget radius $\varepsilon$ is, finally, an explicit economic dial: it trades the
certified coverage level directly against the operating profit the optimizer is permitted
to chase, and exposing that trade-off quantitatively is one of the practical contributions
of the framework.
